%! Complex Numbers and the Library of Parent Functions — Unit 2, Lesson 2.1 — Guided Notes (ANSWER KEY)
% THE in-class document, in the MAIN IDEAS / NOTES shape (LESSON_SHAPE.md §1): the vocabulary
% box, the hook, then ONE two-column guidednotes table. Each row is one idea — a label on the
% left; on the right one or two COMPLETE printed sentences, ONE large pre-drawn display, then a
% two-across grid of numbered problems with real writing room. Problems run 1..18.
% A DENSE DAY: two topics (complex numbers; the parent-function library) plus absolute value
% as a piecewise function, the bridge back to 1.1.
%   I DO   : rows IMAGINARY UNIT and COMPLEX NUMBERS — the teacher works 1, 5 and 7.
%   WE DO  : the rest of those grids (2–4, 6, 8); the PARENT FUNCTIONS gallery (the teacher
%            annotates panels A–C and works 9; the class annotates D–J and works 10); the
%            ABSOLUTE VALUE row (the teacher works 11; the class works 12–14).
%   YOU DO : the last row — problems 15–18, alone, for 12 minutes while the teacher circulates.
% The crux is row 4, ABSOLUTE VALUE: problem 13 tests the claim "-x is negative" at x = -6,
% and problem 14 shows the two readings of |x - 2| disagreeing at x = -1. Problem 18 transfers
% it to |x + 3|, with a negative break point.
%
% PAGE LOCKSTEP with notes_key/: the vocabulary write lines -> \ansline, \blank -> \ans, and
% the answer argument of every \pcell, \writespace and \labelbox filled. Nothing else differs.
% Inside a cell `\\` ENDS THE ROW — break lines with \par. A row cannot break across pages:
% the figure and EACH grid row are separate sub-rows (`\\` then `& ...`, probgrid*).
\documentclass[10pt]{article}
\usepackage{precalculus-article}
\usepackage{precalculus-key}
\newcommand{\TallMath}[1]{$\displaystyle #1\rule[-1.4em]{0pt}{3.2em}$}
% One library card: #1 letter and name, #2 equation, #3 the plot (clipped to the window),
% #4 the domain/range lines. Identical in notes_key/ except for the blanks in #4.
\newcommand{\pfpanel}[4]{%
  \begin{minipage}[t]{0.195\linewidth}\centering
  {\scriptsize\bfseries\color{plum}#1}\par
  {\scriptsize #2}\par\vspace{2pt}
  \begin{tikzpicture}[x=0.25cm, y=0.25cm]
    \draw[linegray, very thin, step=1] (-4,-4) grid (4,4);
    \draw[->, thin] (-4.3,0) -- (4.5,0);
    \draw[->, thin] (0,-4.3) -- (0,4.5);
    \begin{scope}
      \clip (-4,-4) rectangle (4,4);
      #3
    \end{scope}
  \end{tikzpicture}\par\vspace{2pt}
  {\scriptsize #4}
  \end{minipage}}
\newcommand{\kp}[2]{\fill[plum] (#1,#2) circle (1.5pt);}

\begin{document}

\pageheader{Unit 2, Lesson 2.1}{Guided Notes --- Answer Key}

\vspace{0.12in}

\boxguard
\begin{vocabbox}
Fill in each term as we define it together.
\par\vspace{2pt}
\noindent\textbf{\textcolor{plum}{Imaginary unit:}}\\[1pt]
\ansline{The number $i = \sqrt{-1}$, so $i^2 = -1$.}\\[3pt]
\noindent\textbf{\textcolor{plum}{Complex number:}}\\[1pt]
\ansline{A number $a + bi$, with real part $a$ and imaginary part $b$.}\\[3pt]
\noindent\textbf{\textcolor{plum}{Complex conjugate:}}\\[1pt]
\ansline{The partner of $a + bi$ is $a - bi$; their product $a^2 + b^2$ is real.}\\[3pt]
\noindent\textbf{\textcolor{plum}{Parent function:}}\\[1pt]
\ansline{The simplest function of a family, before any shift, stretch, or flip.}\\[3pt]
\noindent\textbf{\textcolor{plum}{Absolute value:}}\\[1pt]
\ansline{The distance from $x$ to 0, never negative: $-x$ for $x < 0$, and $x$ for $x \ge 0$.}
\end{vocabbox}

\vspace{0.1in}

\boxguard
\begin{hookbox}
Type $\sqrt{-9}$ into a calculator and it answers \textsf{ERROR}. Is that because
\emph{no number at all} squares to $-9$ --- or only no number \emph{you have met yet}? What
would a number whose square is negative have to be like?
\ansline{No real number does --- every real square is 0 or positive. So it would have to be a new kind of number, neither positive nor negative. Today we name it $i$; then $3i$ squares to $-9$.}
\end{hookbox}

\small
\begin{guidednotes}
% ── ROW 1: the imaginary unit ────────────────────────────────────────────────
\mainidea[A new number]{Imaginary unit} &
No real number squares to a negative, so mathematicians named one that does. The
\textbf{imaginary unit} is $i = \sqrt{-1}$, which means $i^2 = -1$. For any positive number
$n$, $\sqrt{-n} = i\sqrt{n}$: pull the $i$ out \textbf{first}, then simplify the radical the
usual way.
\\
& \begin{center}
\begin{tikzpicture}[>=stealth]
  \node[font=\normalsize] (a) at (0,1.25) {$i^1 = i$};
  \node[font=\normalsize] (b) at (2.3,0) {$i^2 = -1$};
  \node[font=\normalsize] (c) at (0,-1.25) {$i^3 = $ \labelbox{1.1cm}{$-i$}};
  \node[font=\normalsize] (d) at (-2.4,0) {$i^4 = $ \labelbox{1.1cm}{$1$}};
  \draw[->, very thick, plum] (a) to[bend left=25] node[above right, font=\scriptsize\itshape] {$\times\, i$} (b);
  \draw[->, very thick, plum] (b) to[bend left=25] node[below right, font=\scriptsize\itshape] {$\times\, i$} (c);
  \draw[->, very thick, plum] (c) to[bend left=25] node[below left, font=\scriptsize\itshape] {$\times\, i$} (d);
  \draw[->, very thick, plum] (d) to[bend left=25] node[above left, font=\scriptsize\itshape] {$\times\, i$} (a);
  \node[draw=goldacc, thick, fill=hookbg, rounded corners=2pt, font=\small, align=left,
        text width=4.3cm] at (6.3,0) {\textbf{The cycle:} after $i^4$ the powers repeat. Divide
        the exponent by 4 --- the remainder says where you land.\par\vspace{3pt}
        \textbf{The trap:} $\sqrt{a}\cdot\sqrt{b} = \sqrt{ab}$ only when $a, b \ge 0$.
        Rewrite with $i$ first.};
\end{tikzpicture}
\end{center}
\\
& \notesprompt{Rewrite with $i$ first, then simplify.}
\begin{probgrid}{|Y|Y|}
\pcell{1}{Write $\sqrt{-25}$ and $\sqrt{-12}$ using $i$.}{2.0cm}{$\sqrt{-25} = i\sqrt{25} = 5i$ \par $\sqrt{-12} = i\sqrt{12} = 2i\sqrt{3}$} &
\pcell{2}{Solve $x^2 = -9$. Then solve $x^2 + 16 = 0$.}{2.0cm}{$x = \pm\sqrt{-9} = \pm 3i$. \par $x^2 = -16$, so $x = \pm 4i$.} \\ \hline
\end{probgrid}
\\
& \begin{probgrid}*{|Y|Y|}
\pcell{3}{Find $i^6$ and $i^{11}$. Use the cycle.}{2.0cm}{$6 = 4 + 2$, so $i^6 = i^2 = -1$. \par $11 = 8 + 3$, so $i^{11} = i^3 = -i$.} &
\pcell{4}{A classmate writes $\sqrt{-4}\cdot\sqrt{-9} = \sqrt{36} = 6$. Rewrite each root with $i$ first and find the right answer.}{2.0cm}{$\sqrt{-4}\cdot\sqrt{-9} = (2i)(3i) = 6i^2 = -6$, not 6.} \\ \hline
\end{probgrid}
\\ \hline
% ── ROW 2: complex numbers and their arithmetic ──────────────────────────────
\mainidea[The form $a + bi$]{Complex numbers} &
A \textbf{complex number} has the form $a + bi$, where $a$ and $b$ are real: $a$ is the
\textbf{real part} and $b$ is the \textbf{imaginary part}. Add or subtract by combining real
parts with real parts and imaginary parts with imaginary parts. To multiply:

\vspace{3pt}
\stepnum{1}\ Multiply like binomials (FOIL). \quad \stepnum{2}\ Replace $i^2$ with $-1$.
\quad \stepnum{3}\ Combine like parts; write the answer as $a + bi$.
\\
& \begin{center}
\begin{tikzpicture}[>=stealth]
  \node[font=\Large] (z) at (0,0) {$7 \;-\; 2i$};
  \node (re) at (-1.9,-1.35) {real part: \labelbox{1.0cm}{$7$}};
  \node (im) at (2.1,-1.35) {imaginary part: \labelbox{1.0cm}{$-2$}};
  \draw[->, very thick, plum] (-0.45,-0.3) -- (re.north);
  \draw[->, very thick, plum] (0.55,-0.3) -- (im.north);
  \node[draw=plum, very thick, fill=lilac, rounded corners=3pt, font=\small, align=left,
        text width=4.6cm] at (6.6,-0.6) {The \textbf{complex conjugate} of $a + bi$ is
        $a - bi$ --- flip the sign of the imaginary part. Their product is always real:
        \[ (a + bi)(a - bi) = a^2 + b^2 \]};
\end{tikzpicture}
\end{center}
\\
& \notesprompt{Simplify. Write each answer in the form $a + bi$.}
\begin{probgrid}{|Y|Y|}
\pcell{5}{$(3 + 2i) + (5 - 4i)$}{2.0cm}{$(3 + 5) + (2 - 4)i = 8 - 2i$} &
\pcell{6}{$(7 - i) - (2 - 3i)$. The minus sign reaches both parts.}{2.0cm}{$7 - i - 2 + 3i = 5 + 2i$} \\ \hline
\end{probgrid}
\\
& \begin{probgrid}*{|Y|Y|}
\pcell{7}{$(2 + 3i)(4 - i)$}{2.4cm}{$8 - 2i + 12i - 3i^2$ \par $= 8 + 10i + 3 = 11 + 10i$} &
\pcell{8}{Write the conjugate of $3 + 4i$, then multiply the two. What kind of number is the product?}{2.4cm}{Conjugate: $3 - 4i$. \par $(3 + 4i)(3 - 4i) = 9 - 16i^2 = 9 + 16 = 25$ --- a real number.} \\ \hline
\end{probgrid}
\\ \hline
% ── ROW 3: the library of parent functions ───────────────────────────────────
\mainidea[The library of]{Parent functions} &
A \textbf{parent function} is the simplest function of its family. Every graph in the next
lesson is one of these ten, moved, stretched, or flipped --- so learn each by its
\textbf{shape} and its \textbf{key points} (the dots). Read the domain left to right and the
range bottom to top. Two of them, $\frac{1}{x}$ and $\frac{1}{x^2}$, approach an axis but
never reach it.
\\
& \begin{center}
\pfpanel{A. Constant}{$y = 2$}{%
  \draw[very thick, plum] (-4,2) -- (4,2); \kp{0}{2}}%
  {D: \ans{all reals}\par R: \ans{$y = 2$}}\hfill
\pfpanel{B. Linear}{$y = x$}{%
  \draw[very thick, plum] (-4,-4) -- (4,4); \kp{-1}{-1}\kp{0}{0}\kp{1}{1}}%
  {D: \ans{all reals}\par R: \ans{all reals}}\hfill
\pfpanel{C. Quadratic}{$y = x^2$}{%
  \draw[very thick, plum] plot[domain=-2.05:2.05, samples=50, smooth] (\x,{\x*\x});
  \kp{-1}{1}\kp{0}{0}\kp{1}{1}\kp{-2}{4}\kp{2}{4}}%
  {D: \ans{all reals}\par R: \ans{$y \ge 0$}}\hfill
\pfpanel{D. Cubic}{$y = x^3$}{%
  \draw[very thick, plum] plot[domain=-1.62:1.62, samples=50, smooth] (\x,{\x*\x*\x});
  \kp{-1}{-1}\kp{0}{0}\kp{1}{1}}%
  {D: \ans{all reals}\par R: \ans{all reals}}\hfill
\pfpanel{E. Quartic}{$y = x^4$}{%
  \draw[very thick, plum] plot[domain=-1.43:1.43, samples=50, smooth] (\x,{\x*\x*\x*\x});
  \kp{-1}{1}\kp{0}{0}\kp{1}{1}}%
  {D: \ans{all reals}\par R: \ans{$y \ge 0$}}
\end{center}
\vspace{2pt}
\\
& \begin{center}
\pfpanel{F. Absolute value}{$y = |x|$}{%
  \draw[very thick, plum] (-4,4) -- (0,0) -- (4,4); \kp{-1}{1}\kp{0}{0}\kp{1}{1}}%
  {D: \ans{all reals}\par R: \ans{$y \ge 0$}}\hfill
\pfpanel{G. Square root}{$y = \sqrt{x}$}{%
  \draw[very thick, plum] plot[domain=0:4, samples=60, smooth] (\x,{sqrt(\x)});
  \kp{0}{0}\kp{1}{1}\kp{4}{2}}%
  {D: \ans{$x \ge 0$}\par R: \ans{$y \ge 0$}}\hfill
\pfpanel{H. Cube root}{$y = \sqrt[3]{x}$}{%
  \draw[very thick, plum] plot[domain=-1.62:1.62, samples=50, smooth] ({\x*\x*\x},\x);
  \kp{-1}{-1}\kp{0}{0}\kp{1}{1}}%
  {D: \ans{all reals}\par R: \ans{all reals}}\hfill
\pfpanel{I. Reciprocal}{$y = \frac{1}{x}$}{%
  \draw[very thick, plum] plot[domain=0.25:4, samples=50, smooth] (\x,{1/\x});
  \draw[very thick, plum] plot[domain=-4:-0.25, samples=50, smooth] (\x,{1/\x});
  \kp{1}{1}\kp{-1}{-1}}%
  {D: \ans{$x \ne 0$}\par R: \ans{$y \ne 0$}}\hfill
\pfpanel{J. Reciprocal sq.}{$y = \frac{1}{x^2}$}{%
  \draw[very thick, plum] plot[domain=0.5:4, samples=50, smooth] (\x,{1/(\x*\x)});
  \draw[very thick, plum] plot[domain=-4:-0.5, samples=50, smooth] (\x,{1/(\x*\x)});
  \kp{1}{1}\kp{-1}{1}}%
  {D: \ans{$x \ne 0$}\par R: \ans{$y > 0$}}
\end{center}
\vspace{2pt}
\\
& \notesprompt{Match the description to the library.}
\begin{probgrid}{|Y|Y|}
\pcell{9}{A parent graph passes through $(-8,-2)$, $(-1,-1)$, $(0,0)$, $(1,1)$, and $(8,2)$. Name it and write its equation.}{2.2cm}{The cube root, $y = \sqrt[3]{x}$.} &
\pcell{10}{$y = x^2$ and $y = x^4$ have \emph{even} degree; $y = x^3$ has \emph{odd} degree. Where do the left and right arms of each go? Which would $y = x^5$ look like?}{2.2cm}{Even ($x^2$, $x^4$): both arms point up. Odd ($x^3$): left arm down, right arm up. $x^5$ is odd, so it looks like $x^3$.} \\ \hline
\end{probgrid}
\\ \hline
% ── ROW 4: THE CRUX — absolute value as a piecewise function ─────────────────
\mainidea[A piecewise function]{Absolute value} &
The \textbf{absolute value} $|x|$ is the distance from $x$ to 0, so it is never negative.
Written as a piecewise function,
\[ |x| = \begin{cases} -x & x < 0 \\[2pt] x & x \ge 0 \end{cases} \]
Read $-x$ as \textbf{the opposite of $x$}, not as ``a negative number.'' The rule $-x$ only
ever receives negative inputs, and the opposite of a negative number is positive:
$-(-3) = 3$.
\\
& \begin{center}
\begin{tikzpicture}[x=0.38cm, y=0.38cm, >=stealth]
  \draw[linegray, very thin, step=1] (-5,0) grid (5,5);
  \draw[->, thick] (-5.4,0) -- (5.6,0) node[below right] {\small $x$};
  \draw[->, thick] (0,-0.4) -- (0,5.6) node[above] {\small $|x|$};
  \foreach \x in {-4,-2,2,4} \node[below, font=\scriptsize] at (\x,0) {\x};
  \draw[very thick, redacc] (-5,5) -- (0,0);
  \draw[very thick, plum] (0,0) -- (5,5);
  \fill[plum] (0,0) circle (2.2pt);
  \fill[redacc] (-3,3) circle (2.2pt);
  \node[font=\scriptsize\bfseries\color{redacc}, align=center, left] at (-5.2,2.5) {$-x$ owns\\ $x < 0$};
  \node[font=\scriptsize\bfseries\color{plum}, align=center, right] at (5.4,2.5) {$x$ owns\\ $x \ge 0$};
  \node[font=\scriptsize, right] at (-2.8,3.7) {$(-3,\,3)$};
\end{tikzpicture}
\hspace{0.3cm}
\begin{tikzpicture}[x=0.38cm, y=0.38cm, >=stealth]
  \draw[linegray, very thin, step=1] (-3,0) grid (7,5);
  \draw[->, thick] (-3.4,0) -- (7.6,0) node[below right] {\small $x$};
  \draw[->, thick] (0,-0.4) -- (0,5.6) node[above] {\small $|x - 2|$};
  \foreach \x in {-2,2,4,6} \node[below, font=\scriptsize] at (\x,0) {\x};
  \draw[goldacc, very thick, dashed] (2,0) -- (2,5.3);
  \node[font=\scriptsize\bfseries\color{goldacc}, above right] at (2,5.0) {break at $x = 2$};
  \draw[very thick, redacc] (-3,5) -- (2,0);
  \draw[very thick, plum] (2,0) -- (7,5);
  \fill[plum] (2,0) circle (2.2pt);
\end{tikzpicture}
\end{center}
\vspace{2pt}
$x - 2$ is negative exactly when $x < 2$, so that is where the rule takes the opposite:
\[ |x - 2| = \begin{cases} \text{\labelbox{2.9cm}{$-(x - 2) = 2 - x$}} & x < 2 \\[4pt] \text{\labelbox{2.9cm}{$x - 2$}} & x \ge 2 \end{cases} \]
\\
& \notesprompt{Decide the rule first, then evaluate. Name the rule you used.}
\begin{probgrid}{|Y|Y|}
\pcell{11}{Use the piecewise rule for $|x|$ to find $|-4|$ and $|3|$.}{2.0cm}{$-4 < 0$, so rule $-x$: $|-4| = -(-4) = 4$. \par $3 \ge 0$, so rule $x$: $|3| = 3$.} &
\pcell{12}{Find $-x$ for $x = -5$, $-1$, $0$, and $2$. For which inputs is $-x$ positive?}{2.0cm}{$-x$ is $5$, $1$, $0$, $-2$. It is positive exactly when $x < 0$ --- the inputs the rule $-x$ owns.} \\ \hline
\end{probgrid}
\\
& \begin{probgrid}*{|Y|Y|}
\pcell{13}{A classmate says: ``$|x| = -x$ for $x < 0$ can't be right, because $-x$ is negative.'' Test the claim at $x = -6$. Who is right, and why?}{2.8cm}{$-x = -(-6) = 6$, positive, and $|-6| = 6$. The rule is right: $-x$ means \emph{the opposite of} $x$, and the opposite of a negative is positive.} &
\pcell{14}{Using the $|x - 2|$ rules above, find $|x - 2|$ at $x = -1$ and at $x = 5$. What would $x - 2$ alone give at $x = -1$, and why can't that be right?}{2.8cm}{$-1 < 2$: $2 - (-1) = 3$. \ $5 \ge 2$: $5 - 2 = 3$. \par $x - 2$ at $x = -1$ gives $-3$ --- but an absolute value is a distance, never negative.} \\ \hline
\end{probgrid}
\\ \hline
% ── LAST ROW: YOU DO — alone, while the teacher circulates (12 min) ──────────
\mainidea[You do, alone]{All three} &
Complex arithmetic, one graph from the library, and an absolute value written as a piecewise
function. The graph for problem 17:
\\
& \begin{center}
\begin{tikzpicture}[x=0.45cm, y=0.45cm, >=stealth]
  \draw[linegray, very thin, step=1] (-4,-1) grid (4,5);
  \draw[->, thick] (-4.4,0) -- (4.6,0) node[below right] {\small $x$};
  \draw[->, thick] (0,-1.4) -- (0,5.6) node[above] {\small $y$};
  \foreach \x in {-3,-2,-1,1,2,3} \node[below, font=\scriptsize] at (\x,0) {\x};
  \foreach \y in {1,2,3,4} \node[left, font=\scriptsize] at (0,\y) {\y};
  \begin{scope}
    \clip (-4,-1) rectangle (4,5);
    \draw[very thick, plum] plot[domain=0.45:4, samples=60, smooth] (\x,{1/(\x*\x)});
    \draw[very thick, plum] plot[domain=-4:-0.45, samples=60, smooth] (\x,{1/(\x*\x)});
  \end{scope}
  \fill[plum] (1,1) circle (2.2pt);
  \fill[plum] (-1,1) circle (2.2pt);
  \fill[plum] (0.5,4) circle (2.2pt);
  \fill[plum] (-0.5,4) circle (2.2pt);
\end{tikzpicture}
\end{center}
\\
& \notesprompt{On your own. Show your work.}
\begin{probgrid}{|Y|Y|}
\pcell{15}{Simplify $(4 - 3i)(2 + i)$. Write the answer as $a + bi$.}{3.2cm}{$8 + 4i - 6i - 3i^2$ \par $= 8 - 2i + 3 = 11 - 2i$} &
\pcell{16}{Solve $x^2 + 36 = 0$. Then find $\sqrt{-9}\cdot\sqrt{-16}$.}{3.2cm}{$x^2 = -36$, so $x = \pm 6i$. \par $(3i)(4i) = 12i^2 = -12$} \\ \hline
\end{probgrid}
\\
& \begin{probgrid}*{|Y|Y|}
\pcell{17}{The graph above is a parent function. Name it, write its equation, and state its domain and range.}{3.2cm}{Reciprocal squared, $y = \frac{1}{x^2}$. \par Domain $x \ne 0$; range $y > 0$.} &
\pcell{18}{Write $|x + 3|$ as a piecewise function. (Where does $x + 3$ change sign?) Then find its value at $x = -5$ and $x = 1$, naming the rule each time.}{3.4cm}{$|x + 3| = -(x + 3)$ for $x < -3$, and $x + 3$ for $x \ge -3$. \par $-5 < -3$: $-(-5 + 3) = 2$. \par $1 \ge -3$: $1 + 3 = 4$.} \\ \hline
\end{probgrid}
\\ \hline
\end{guidednotes}

% ── THE NOTES END AT THE YOU DO ROW; AP Practice and the Homework follow ──────

\end{document}
